\subsection{Dual of a direct sum and of a product}

For every \(i \in I\) , let \((\mathbf{E}_{i}, \mathbf{F}_{i})\) be a pair of vector spaces, set in duality by a bilinear form \(B_{i}\) . We put \(E = \prod_{i \in I} E_{i}\) and \(F = \oplus_{i \in I} F_{i}\) , and we identify each \(F_{i}\) with a subspace of F. We put E and F in duality by means of the bilinear form

\[
\mathrm{B} (x, y) = \sum_ {i \in \mathrm{I}} \mathrm{B} _ {i} \left(x _ {i}, y _ {i}\right) \quad \text { for } \quad x = \left(x _ {i}\right) \quad \text { and } \quad y = \left(y _ {i}\right)\tag{8}
\]

(the family \((\mathbf{B}_i(x_i, y_i))_{i \in I}\) has finite support).

We recall (II, p.~50, prop. 8) that the weak topology \(\sigma(\mathrm{E},\mathrm{F})\) is the product of the weak topologies \(\sigma(\mathrm{E}_i,\mathrm{F}_i)\) .

Lemma 3. --- (i) For every \(i \in \mathbf{I}\) , let \(\mathfrak{S}_i\) be a family of subsets of \(\mathrm{F}_i\) , which is bounded for \(\sigma(\mathrm{F}_i, \mathrm{E}_i)\) ; put \(\mathfrak{S} = \bigcup_{i \in \mathbf{I}} \mathfrak{S}_i\) . Then the \(\mathfrak{S}\) -topology on \(\mathbf{E}\) is the product of the \(\mathfrak{S}_i\) -topologies on the \(\mathrm{E}_i\) .

\begin{enumerate}
\def\labelenumi{(\roman{enumi})}
\setcounter{enumi}{1}
\tightlist
\item
  For every \(i \in \mathbf{I}\) , let \(\mathfrak{J}_i\) be an adapted bornology on the space \(\mathrm{E}_i\) endowed with the weak topology \(\sigma(\mathrm{E}_i, \mathrm{F}_i)\) , none equal to \(\{\varnothing\}\) . Let \(\mathfrak{J}\) be the family of subsets \(\mathbf{A}\) of \(\mathrm{E} = \prod_{i \in \mathbf{I}} \mathrm{E}_i\) such that \(\mathrm{pr}_i(\mathbf{A}) \in \mathfrak{J}_i\) for all \(i \in \mathbf{I}\) . Then the \(\mathfrak{J}\) -topology on \(\mathbf{F}\) is the direct sum of the \(\mathfrak{J}_i\) -topologies on the \(\mathrm{F}_i\) .
\end{enumerate}

Let \(\mathcal{T}\) be the product of the \(\mathfrak{S}_i\) -topologies. The sets of the form

\[
\mathbf {A} = \prod_ {i \in \mathbf {J}} \mathbf {A} _ {i} ^ {\circ} \times \prod_ {i \in \mathbf {I} - \mathbf {J}} \mathbf {E} _ {i}
\]

where \(\mathbf{J} \subset \mathbf{I}\) is finite and \(\mathbf{A}_i \in \mathfrak{S}_i\) for all \(i \in \mathbf{J}\) , form a fundamental system of neighbourhoods of 0 for \(\mathcal{T}\) . We have \(\mathbf{A} = (\bigcup_{i \in \mathbf{J}} \mathbf{A}_i)^\circ\) , hence \(\mathcal{T}\) is identical with the \(\mathfrak{S}\) -topology. This proves (i).

We assign the \(\mathfrak{J}\) -topology to \(F\) and the \(\mathfrak{J}_i\) -topology to each \(F_i\) . For every subset \(A\) of \(E\) , we have \(F_i \cap A^\circ = \text{pr}_i(A)^\circ\) , hence the injection from \(F_i\) into \(F\) is continuous. Let \(q\) be a semi-norm on \(F\) ; we assume that the restriction \(q_i\) of \(q\) to \(F_i\) is continuous for all \(i \in I\) . Then we can find non-empty subsets \(A_i \in \mathfrak{J}_i\) such that we have

\[
q _ {i} \left(y _ {i}\right) \leqslant \sup _ {x _ {i} \in \mathrm{A} _ {i}} \left| \mathrm{B} _ {i} \left(x _ {i}, y _ {i}\right) \right| \quad \left(y _ {i} \in \mathrm{F} _ {i}\right).\tag{9}
\]

Put \(\mathrm{A} = \prod_{i\in \mathrm{I}}\mathrm{A}_i\) ; then \(\mathrm{A}\in \mathfrak{J}\) . For \(y = (y_i)_{i\in \mathrm{I}}\) in \(\mathrm{F}\) , we have

\[
q (y) \leqslant \sum_ {i \in \mathbf {I}} q _ {i} (y _ {i}) \leqslant \sum_ {i \in \mathbf {I}} \sup _ {x _ {i} \in \mathrm{A} _ {i}} | \mathrm{B} _ {i} (x _ {i}, y _ {i}) | = \sup _ {x \in \mathrm{A}} | \mathrm{B} (x, y) |,
\]

where the last equality follows from (8) since the family \((y_{i})_{i\in\mathbf{I}}\) has finite support and the \(A_{i}\) are non-empty and can be assumed balanced (GT, IV, § 5, No.~7, cor. 2 to prop. 12). This inequality proves that q is continuous on F, and hence (ii).

PROPOSITION 12. --- The topology \(\beta(\mathrm{F},\mathrm{E})\) is the direct sum of the topologies \(\beta(F_i,E_i)\) . The topology \(\beta(\mathrm{E},\mathrm{F})\) is the product of the topologies \(\beta(\mathrm{E}_i,\mathrm{F}_i)\) .

We shall apply lemma 3 taking for \(\mathfrak{S}_i\) the family of all subsets of \(F_i\) which are bounded for \(\sigma(F_i, E_i)\) and for \(\mathfrak{J}_i\) the family of all subsets of \(E_i\) which are bounded for \(\sigma(E_i, F_i)\) .

By cor. 2 of III, p.~4, \(\mathfrak{J}\) is the family of all subsets of \(E_{i}\) which are bounded for the product topology of the \(\sigma(\mathrm{E}_{i}, \mathrm{F}_{i})\) , which is identical with \(\sigma(\mathrm{E}, \mathrm{F})\) . Hence our assertion on \(\beta(\mathrm{F}, \mathrm{E})\) follows.

We endow \(\mathrm{F} = \bigoplus_{i\in \mathrm{I}}\mathrm{F}_i\) with the topology \(\mathcal{T}\) which is the direct sum of the topologies \(\sigma (\mathrm{F}_i,\mathrm{E}_i)\) . Then the dual of \(\mathrm{F}\) consists of the linear forms \(y\mapsto \mathrm{B}(x,y)\) where \(x\) runs through \(\mathrm{E}\) (II, p.~30, prop. 6). By prop. 1 of IV, p.~1, the topologies \(\mathcal{T}\) and \(\sigma (\mathrm{F},\mathrm{E})\) have the same bounded sets. Assume first that the topologies \(\sigma (\mathrm{F}_i,\mathrm{E}_i)\) are Hausdorff. By prop. 5 of III, p.~5, these sets are contained in a subset of the form \(\sum_{i\in \mathrm{J}}\mathrm{B}_i\) with \(J\subset I\) finite and \(B_{i}\) bounded in \(F_{i}\) (for \(\sigma (\mathrm{F}_i,\mathrm{E}_i)\) ) for all \(i\in J\) . Since \(\sum_{i\in J}\mathrm{B}_{i}\) is contained in the convex envelope of \(\bigcup_{i\in J}n\mathrm{B}_{i}\) , where \(n = \operatorname {Card}(J)\) , we can apply lemma 3, to prove the assertion on \(\beta (\mathrm{E},\mathrm{F})\) in this case.

For the general case, let \(\mathbf{N}_i\) be the intersection of all neighbourhoods of 0 for \(\sigma (\mathrm{F}_i,\mathrm{E}_i)\) , and let \(\mathbf{N} = \sum_{i\in \mathbf{I}}\mathbf{N}_i\) ; then \(\mathbf{F} / \mathbf{N}\) is the topological direct sum of the \(\mathrm{F}_i / \mathrm{N}_i\) (II, p.~31, prop. 8); we deduce from this that every subset of \(\mathbf{F}\) which is bounded for \(\mathcal{T}\) is contained in a set of the form \(\mathbf{N} + \sum_{i\in \mathbf{J}}\mathbf{B}_i\) with \(\mathrm{J}\subset \mathrm{I}\) finite and \(\mathbf{B}_i\) bounded in \(\mathbf{F}_i\) for all \(i\in \mathbf{J}\) (III, p.~2, Remark 3); since the polar of this set in \(\mathbf{E}\) is the same as that of \(\sum_{i\in \mathbf{J}}\mathbf{B}_i\) , the result follows as above.

PROPOSITION 13. --- The Mackey topology \(\tau(\mathrm{F},\mathrm{E})\) is the direct sum of the Mackey topologies \(\tau(\mathrm{F}_i,\mathrm{E}_i)\) . The topology \(\tau(\mathrm{E},\mathrm{F})\) is the product of the topologies \(\tau(\mathrm{E}_i,\mathrm{F}_i)\) .

The assertion on \(\tau(\mathrm{F},\mathrm{E})\) follows from lemma 3 (ii) and the following property: for a closed, convex and balanced subset of \(F^{*} = \prod_{i\in I}F_{i}^{*}\) to be compact for \(\sigma(F^{*},F)\) , it is necessary and sufficient that its projection on each \(\mathrm{F}_i^*\) is compact for \(\sigma(\mathrm{F}_i^*, \mathrm{F}_i)\) .

To prove the assertion on \(\tau(\mathrm{E},\mathrm{F})\) , assume first that the topologies \(\sigma(\mathrm{F}_i,\mathrm{E}_i)\) are Hausdorff, it is enough (lemma 3 (i)) to prove that every subset A of F which is convex, balanced and compact for \(\sigma(\mathrm{F},\mathrm{E})\) is contained in a set of the form \(\sum_{i\in\mathbf{J}}\mathrm{A}_i\) where

\(J \subset I\) is finite and where \(A_{i}\) is convex, balanced and compact for \(\sigma(F_{i}, E_{i})\) . But such a subset is bounded for \(\sigma(F, E)\) . By the proof of prop. 12, there exists a finite subset \(J\) of \(I\) such that \(A \subset \sum_{i \in J} F_i\) , and it is enough to take for \(A_i\) the projection of \(A\) on \(F_i\) .

In the general case, with the same notations as in the proof of prop. 12, we have \(\tau(\mathrm{E}_{i},\mathrm{F}_{i})=\tau(\mathrm{E}_{i},\mathrm{F}_{i}/\mathrm{N}_{i})\) and \(\tau(\mathrm{E},\mathrm{F})=\tau(\mathrm{E},\mathrm{F}/\mathrm{N})\) (IV, p.~2) and since F/N is the topological direct sum of the \(F_{i}/N_{i}\) , we have reduced to the preceding case.

Q.E.D.

For the remainder of this paragraph, we assume that \((\mathrm{E}_i)_{i\in \mathbf{I}}\) is a family of locally convex spaces. Let S denote the topological direct sum of the \(\mathrm{E}_i\) and P, their product. We define a linear mapping \(\theta :\mathrm{S}'\to \prod_{i\in \mathbf{I}}\mathrm{E}_i'\) , said to be canonical, by

\[
\theta (x ^ {\prime}) = (x ^ {\prime} | \mathrm{E} _ {i}) _ {i \in \mathrm{I}} (x ^ {\prime} \in \mathrm{S} ^ {\prime})\tag{10}
\]

(where \(S'\) denotes the dual of \(S\) , and \(E_i'\) that of \(E_i\) ).

PROPOSITION 14. --- (i) The mapping \(\theta\) is an isomorphism from the strong (resp. weak) dual of \(S = \oplus E_i\) onto the product of the strong (resp. weak) duals of the \(E_i\) :

\begin{enumerate}
\def\labelenumi{(\roman{enumi})}
\setcounter{enumi}{1}
\item
  For a subset A of \(\mathbf{S}'\) to be equicontinuous, it is necessary and sufficient that the projection of \(\theta (\mathbf{A})\) onto \(\mathrm{E}_i'\) be equicontinuous for all \(i\in \mathbf{I}\) .
\item
  The Mackey topology \(\tau(\mathbf{S},\mathbf{S}')\) is the direct sum of the Mackey topologies \(\tau(\mathbf{E}_i,\mathbf{E}_i')\) .
\item
  The topology \(\beta (\mathbf{S},\mathbf{S}^{\prime})\) is the direct sum of the topologies \(\beta (\mathrm{E}_i,\mathrm{E}_i^{\prime})\).
\end{enumerate}

That \(\theta\) is bijective follows immediately from the definition of a topological direct sum (II, p.~30, prop. 6). Assertion (i) then follows from prop. 12 of IV, p.~12, for the strong topologies, and from prop. 8 of II, p.~50, for the weak topologies. Similarly (iii) follows from prop. 13 (IV, p.~12) and (iv) from prop. 12 (IV, p.~12).

To prove (ii), let A be a subset of S'. Put

\[
q (x) = \sup _ {x ^ {\prime} \in \mathbf {A}} | \langle x, x ^ {\prime} \rangle | \quad \text { for } \quad x \in \mathrm{S};\tag{11}
\]

let \(q_{i}\) denote the restriction of q to \(E_{i}\) , whence

\[
q _ {i} (x _ {i}) = \sup _ {x _ {i} ^ {\prime} \in \mathbf {A} _ {i}} | \langle x _ {i}, x _ {i} ^ {\prime} \rangle | \quad \text { for } \quad x _ {i} \in \mathrm{E} _ {i},\tag{12}
\]

where \(A_{i}\) denotes the projection of \(\theta(\mathbf{A})\) on \(E_{i}^{\prime}\) . For A to be equicontinuous, it is necessary and sufficient that q is finite (that is, that each \(q_{i}\) is finite) and continuous. In view of the characterization of continuous semi-norms on a topological direct sum (II, p.~27, prop. 5), this is the same as assuming that each \(q_{i}\) is continuous, or in fact, that each set \(A_{i}\) is equicontinuous. Q.E.D.

Let \(\phi\) be the linear mapping, said to be canonical, from \(\oplus E_i'\) into the dual \(P'\)

of \(\mathbf{P} = \prod_{i\in \mathbf{I}}\mathrm{E}_i\) , defined by the formula

\[
\langle x, \phi (x ^ {\prime}) \rangle = \sum_ {i \in \mathbf {I}} \langle x _ {i}, x _ {i} ^ {\prime} \rangle\tag{13}
\]

for \(x = (x_{i})\) in \(\mathrm{P}\) and \(x^{\prime} = (x_{i}^{\prime})\) in \(\bigoplus_{i\in \mathbf{I}}\mathrm{E}_{i}^{\prime}\) .

PROPOSITION 15. --- (i) The map \(\phi\) is an isomorphism from the topological direct sum of the strong duals of the \(E_i\) onto the strong dual of \(P = \prod_{i \in I} E_i\) .

\begin{enumerate}
\def\labelenumi{(\roman{enumi})}
\setcounter{enumi}{1}
\item
  For a subset \(\mathbf{A}\) of \(\mathbf{P}'\) to be equicontinuous, it is necessary and sufficient that it is contained in a finite sum \(\sum_{i\in \mathbf{J}}\phi (\mathbf{A}_i)\) , where \(\mathbf{J}\subset \mathbf{I}\) is finite and where \(\mathbf{A}_i\) is equicontinuous in \(\mathbf{E}_i'\) for all \(i\in \mathbf{J}\) .
\item
  The Mackey topology \(\tau (\mathbf{P},\mathbf{P}^{\prime})\) is the product of the topologies \(\tau (\mathrm{E}_i,\mathrm{E}_i')\).
\item
  The topology \(\beta (\mathbf{P},\mathbf{P}^{\prime})\) is the product of the topologies \(\beta (\mathrm{E}_i,\mathrm{E}_i^{\prime})\).
\end{enumerate}

It is immediate that \(\phi\) is injective. A fundamental system of neighbourhoods of 0 in \(\mathbf{P}\) consists of sets of the form \(\mathrm{V} = \prod_{i\in \mathbf{J}}\mathrm{V}_i\times \prod_{i\in \mathbf{I} - \mathbf{J}}\mathrm{E}_i\) , where \(\mathrm{J}\subset \mathrm{I}\) is finite and \(\mathrm{V}_i\) is a neighbourhood of 0 in \(\mathrm{E}_i\) for \(i\) in \(\mathrm{J}\) . The polar of \(\mathrm{V}\) in \(\mathbf{P}'\) is equal to \(\sum_{i\in \mathbf{J}}\phi (\mathrm{V}_i^\circ)\) . This proves the surjectivity of \(\phi\) and also assertion (ii).

Assertions (i) and (iv) follow from prop. 12 (IV, p.~12) and (iii) from prop. 13 (IV, p.~12).

COROLLARY. --- Every product of barrelled spaces is barrelled.

A locally convex space E is barrelled if and only if the initial topology is identical with \(\beta(\mathrm{E}, \mathrm{E}')\) (IV, p.~4, Remark 3). Hence it is enough to apply prop. 15 (iv).

